% !TEX program = xelatex
\documentclass{article}
\usepackage{geometry}\geometry{margin=23mm}
\usepackage{fontspec}\setmainfont{Times New Roman}
\usepackage[portuguese]{babel}
\begin{document}
\title{A tábua de medição}\author{}\date{14 de maio de 2026}\maketitle
\begin{abstract}Duas marcas indicam a posição inicial da tábua.\end{abstract}
\tableofcontents
\section{As marcas}\begin{quote}«Ler primeiro a marca inferior.»\end{quote}
\begin{figure}[ht]\centering\fbox{\rule{0pt}{8mm}\rule{33mm}{0pt}}\caption{Posição da tábua}\end{figure}
\section{Altura}
\begin{table}[ht]\centering\begin{tabular}{cr}Marca&Altura (mm)\\\hline 1&12\\2&18\end{tabular}\caption{Altura das marcas}\end{table}
\begin{thebibliography}{1}\bibitem{record}\emph{Caderno de medições}, 2025, \chaptername~2.\end{thebibliography}
\end{document}
